\chapter*{Abstract}
\addcontentsline{toc}{chapter}{Abstract}

Nursing homes near the coast have to move their residents before a storm surge
arrives. This needs several organizations to work together, and the plan often
changes on the way. Simulations of this task are usually scripts that cover
only the disruptions their authors expected. Agents built from large language
models (LLMs) could, in principle, also handle disruptions that nobody scripted.

We design a database-centric agentic architecture and compare it with a
scripted baseline that follows the German DV~100 doctrine. The case is the
evacuation of the Augustinum nursing home in Hamburg with 100 residents. Ten
agents represent people at eight organizations. Six of them reason with an
LLM, and all coordinate through one PostgreSQL database. Both systems write to
one shared event log, and one deterministic scorer reads it. We compared them
on seven scenarios under a pre-registered protocol.

Both pre-registered hypotheses failed. The baseline sheltered every resident
in all 50 runs with a disruption it was not written for, but it did not adapt.
It kept its original orders. A controller that ignores every disruption scores
the same, because no disruption changed what the scorer counts. The agentic
system did react, and its agents read the disruptions correctly. Yet the best
configuration sheltered only 89.4\% of the residents, and in every scenario
with such a disruption it did worse than ignoring the disruption. The lost
residents were never dispatched, so the loss came in whole bus-loads. The more
often a run revised its plan, the fewer residents it delivered. Coordination
itself worked, with 93.8\% of the requests between agents answered. We
therefore see a failure of commitment rather than of comprehension. Since no
scenario required adaptation, these results show who won here, not whether
adaptation can ever pay off.

From this we derive design considerations. An agent should not cancel a bus run
once it has started, every delivery should be confirmed, and test scenarios
should be ones where ignoring the disruption costs residents. As a next step,
we propose to rerun the same scenarios with agents that may revise their plan
only a few times.

\chapter*{Acknowledgements}
\addcontentsline{toc}{chapter}{Acknowledgements}

I would like to thank Prof. Dr.-Ing. habil. Kai-Uwe Sattler for the chance to
write this thesis in the Databases and Information Systems Group, and for the
guidance along the way. I am also grateful to M.\,Sc. Saba Zamankhani for the
supervision, the time given to this work and the helpful feedback on the design
and the evaluation.

This thesis was written at Absolute Software GmbH in Hamburg, as part of the
company's work in the RESCUE-MATE project. RESCUE-MATE is a research project
funded under the German programme Research for Civil Security. It studies how
rescue forces can be supported in complex crisis situations, and its main
scenario is a severe storm surge in Hamburg that requires large evacuations
near the Elbe. This scenario gave my thesis its case. I thank Absolute
Software GmbH for the opportunity to work on this topic within the project,
and for the support I received during the whole thesis. My special thanks go
to my company supervisor, Mr. Rüdiger Höfert, for the trust, the
encouragement and the many helpful discussions along the way.
